%% kapa: exam=1 round=34 year=2023 period=1 form=A subject=real_estate subject_ko=부동산학원론 seq=38 number=118
%% source: Q-Net 감정평가사 기출문제 (cst00309) artlSeq=5211789 file=raw/1008/5211789/2023년도 제34회 감정평가사 1차시험 1교시.pdf
%% conversion: kordoc 4.24.51 -> tools/gen1.py (auto); review=proofread-sonnet; flags=none
\begin{question}{118}
다음 조건을 가진 부동산을 통해 산출한 내용으로 옳지 않은 것은? (단, 주어진 조건에 한함)
\begin{framedtext}
○ 가능총소득(PGI): 연 150,000,000원 \par
○ 공실손실상당액·대손충당금: 가능총소득의 10 \% \par
○ 운영경비(OE): 유효총소득의 30 \% \par
○ 대출원리금 상환액: 연 40,000,000원 \par
○ 가격구성비: 토지 40 \%, 건물 60 \% \par
○ 토지환원이율: 연 3 \%, 건물환원이율: 연 5 \% \par
\end{framedtext}
\begin{choices}
  \item 운영경비는 40,500,000원이다.
  \item 종합환원이율은 연 4.2 \%이다.
  \item 순영업소득(NOI)은 94,500,000원이다.
  \item 유효총소득(EGI)은 135,000,000원이다.
  \item 세전현금흐름(BTCF)은 53,500,000원이다.
\end{choices}
\end{question}
\answer{5}
